\documentclass[10pt,a4paper]{amsart}
\usepackage{amsmath,amssymb,amsthm,graphicx}
% Retained preamble for the Stacks Project chapter snapshot in this directory.
% The chapter's own preamble.tex is a copy of the shared Stacks Project preamble; the block
% below is sliced MECHANICALLY, line for line, from it. Nothing was added, removed or
% redefined; the wrapper class, the xr-hyper/hyperref wiring and the externaldocument list
% are omitted because the assembler supplies the document class.
%
% lines taken from preamble.tex: 3, 5, 9, 11, 12, 156, 157, 158, ... (51 lines total)
%
% package coverage restored from the chapter's own preamble (2026-10-13): xy, xy-2cell

\usepackage[all]{xy}\usepackage{verbatim}
\xyoption{2cell}
\UseAllTwocells

\newenvironment{reference}{\comment}{\endcomment}
\newenvironment{slogan}{\comment}{\endcomment}
\newenvironment{history}{\comment}{\endcomment}
\theoremstyle{plain}
\newtheorem{theorem}[subsection]{Theorem}
\newtheorem{proposition}[subsection]{Proposition}
\newtheorem{lemma}[subsection]{Lemma}
\theoremstyle{definition}
\newtheorem{definition}[subsection]{Definition}
\newtheorem{example}[subsection]{Example}
\newtheorem{exercise}[subsection]{Exercise}
\newtheorem{situation}[subsection]{Situation}
\theoremstyle{remark}
\newtheorem{remark}[subsection]{Remark}
\newtheorem{remarks}[subsection]{Remarks}
\numberwithin{equation}{subsection}
\def\lim{\mathop{\mathrm{lim}}\nolimits}
\def\colim{\mathop{\mathrm{colim}}\nolimits}
\def\Spec{\mathop{\mathrm{Spec}}}
\def\Hom{\mathop{\mathrm{Hom}}\nolimits}
\def\Ext{\mathop{\mathrm{Ext}}\nolimits}
\def\SheafHom{\mathop{\mathcal{H}\!\mathit{om}}\nolimits}
\def\SheafExt{\mathop{\mathcal{E}\!\mathit{xt}}\nolimits}
\def\Sch{\mathit{Sch}}
\def\Mor{\mathop{\mathrm{Mor}}\nolimits}
\def\Ob{\mathop{\mathrm{Ob}}\nolimits}
\def\Sh{\mathop{\mathit{Sh}}\nolimits}
\def\NL{\mathop{N\!L}\nolimits}
\def\CH{\mathop{\mathrm{CH}}\nolimits}
\def\proetale{{pro\text{-}\acute{e}tale}}
\def\etale{{\acute{e}tale}}
\def\QCoh{\mathit{QCoh}}
\def\Ker{\mathop{\mathrm{Ker}}}
\def\Im{\mathop{\mathrm{Im}}}
\def\Coker{\mathop{\mathrm{Coker}}}
\def\Coim{\mathop{\mathrm{Coim}}}
\DeclareMathSymbol{\boxtimes}{\mathbin}{AMSa}{"02}
\def\QCohstack{\mathcal{QC}\!\mathit{oh}}
\def\Cohstack{\mathcal{C}\!\mathit{oh}}
\def\Spacesstack{\mathcal{S}\!\mathit{paces}}
\def\Quotfunctor{\mathrm{Quot}}
\def\Hilbfunctor{\mathrm{Hilb}}
\def\Curvesstack{\mathcal{C}\!\mathit{urves}}
\def\Polarizedstack{\mathcal{P}\!\mathit{olarized}}
\def\Complexesstack{\mathcal{C}\!\mathit{omplexes}}
\def\Pic{\mathop{\mathrm{Pic}}\nolimits}
\def\Picardstack{\mathcal{P}\!\mathit{ic}}
\def\Picardfunctor{\mathrm{Pic}}
\def\Deformationcategory{\mathcal{D}\!\mathit{ef}}
% lane-s3 patch-sd-envs.tex - DO NOT REMOVE LINES 16-21 (the executable patch).
% Contract (Chief-confirmed 2026-10-14):
%  * Lines 16-18: reference/slogan/history are declared as \comment containers in
%    every retained-preamble of this lane; the patch empties them so navigation
%    pointer text inside statements typesets without "Environment undefined".
%  * Line 21: \cite is globally no-op'd because the corpus rule is "cards carry
%    no bibliography": there is no bibliography backend at assembly time, and
%    every patched card's statement+proof span was scanned zero-cite before
%    patching (2026-10-14, lane-s3) - the only \cite occurrences live inside
%    the emptied reference environments (e.g. \cite{EGA} pointers). Removing
%    line 21 breaks compilation of every card whose statement carries a
%    reference environment with a citation (co01 wave: Citation `EGA' undefined).
%  * The patch is baked into task.json at card-build time (preamble_span
%    end_line must be the last line of this file - the builder computes it
%    dynamically since 2026-10-14); cards without patchenv never load it.
\renewenvironment{reference}{}{}
\renewenvironment{slogan}{}{}
\renewenvironment{history}{}{}
% empty out citations inside navigation-pointer environments (no bibliography
% backend; spans themselves are cite-free - scanned 2026-10-14, lane-s3)
\renewcommand{\cite}[2][]{}
\usepackage{cleveref}
\crefname{originalproposition}{original proposition}{original propositions}
\Crefname{originalproposition}{Original proposition}{Original propositions}
\makeatletter
\newlabel{descent-lemma-quasi-affine}{{0247}{source}{}{sourcepointer-1}{}}
\newlabel{descent-lemma-quasi-affine@cref}{{[originalproposition][0][]0247}{[1][0][]source}{}{}{}}
\makeatother
\crefname{originalproposition}{original proposition}{original propositions}
\Crefname{originalproposition}{Original proposition}{Original propositions}
\makeatletter
\newlabel{more-groupoids-lemma-ind-quasi-affine}{{0APK}{source}{}{sourcepointer-2}{}}
\newlabel{more-groupoids-lemma-ind-quasi-affine@cref}{{[originalproposition][0][]0APK}{[1][0][]source}{}{}{}}
\makeatother
\crefname{originalproposition}{original proposition}{original propositions}
\Crefname{originalproposition}{Original proposition}{Original propositions}
\makeatletter
\newlabel{spaces-properties-proposition-sheaf-fpqc}{{0APL}{source}{}{sourcepointer-3}{}}
\newlabel{spaces-properties-proposition-sheaf-fpqc@cref}{{[originalproposition][0][]0APL}{[1][0][]source}{}{}{}}
\makeatother
\crefname{originalproposition}{original proposition}{original propositions}
\Crefname{originalproposition}{Original proposition}{Original propositions}
\makeatletter
\newlabel{stacks-section-descent-data}{{02ZC}{source}{}{sourcepointer-4}{}}
\newlabel{stacks-section-descent-data@cref}{{[originalproposition][0][]02ZC}{[1][0][]source}{}{}{}}
\makeatother
\title{Stacks with quasi-affine diagonal satisfy fpqc descent}
\author{The Stacks Project Authors}
\date{stacks-project:stacks-more-morphisms.tex:local-snapshot-20261012; GNU FDL 1.2 or later}
\begin{document}
\maketitle
\noindent\textit{Editorial note.} Statement span is lines 4320-4329 of sources/stacks-stacks-more-morphisms/pinned-original-source.tex; proof span is lines 4331-4457, the complete author proof (125 body lines).\par\medskip
\noindent\textit{Editorial note.} Navigation stubs: 4 cited-but-not-included labels are supplied as non-load-bearing pointers with their pinned tags where the snapshot tag index has them.\par\medskip
\noindent\textit{Editorial note.} extra\_packages is empty on purpose: the retained preamble already loads xy with options and a second declaration clashes in the assembler wrapper.\par\medskip
\noindent\textit{Editorial note.} Text is retained verbatim, including the author's own typos and placeholders.\par\medskip
\noindent\textit{Editorial note.} Queue index 4 of queue\_stacks-stacks-more-morphisms.json (pb=4561, stmt 4320-4329, proof 4331-4457).\par\medskip
\noindent\textit{Editorial note.} stacks-stacks-more-morphisms wave smm01 (2026-10-14): fresh adjudication accept per lane-s3/smm\_judgments.md; pool per Chief 10:5x open-pool ruling (net 59 with-proof rows, bloodline lane-s3 by anchor audit, v2 true-key namespace stacks:stacks-more-morphisms.tex, row lemma-deform-property-fp-over-ft avoided on cross-chapter key collision); 60 with-proof rows judged 45 accept / 14 reject / 1 avoid; patch-sd-envs contract spans scanned.\par\medskip
\noindent\textit{Editorial note.} Statement carries a reference/slogan/history environment (patchenv -> patch-sd-envs.tex).\par\medskip
\noindent\textit{Allowed standard prerequisites (editorial).}\begin{itemize}
\item Explicit prerequisite boundary. The cited inputs are prerequisites; the argument in the retained span is not replaced by any cited result.
\end{itemize}\medskip
\ifdefined\hypertarget\hypertarget{sourcepointer-1}{}\fi
\noindent\textit{Source reference.} Original proposition 0247: sources/stacks-stacks-more-morphisms/pinned-original-source.tex, source lines 4457--4457. Cited by the retained statement or proof but not part of the collected spans; supplied only as a navigation pointer. The argument does not depend on its content.\par\medskip
\ifdefined\hypertarget\hypertarget{sourcepointer-2}{}\fi
\noindent\textit{Source reference.} Original proposition 0APK: sources/stacks-stacks-more-morphisms/pinned-original-source.tex, source lines 4457--4457. Cited by the retained statement or proof but not part of the collected spans; supplied only as a navigation pointer. The argument does not depend on its content.\par\medskip
\ifdefined\hypertarget\hypertarget{sourcepointer-3}{}\fi
\noindent\textit{Source reference.} Original proposition 0APL: sources/stacks-stacks-more-morphisms/pinned-original-source.tex, source lines 4457--4457. Cited by the retained statement or proof but not part of the collected spans; supplied only as a navigation pointer. The argument does not depend on its content.\par\medskip
\ifdefined\hypertarget\hypertarget{sourcepointer-4}{}\fi
\noindent\textit{Source reference.} Original proposition 02ZC: sources/stacks-stacks-more-morphisms/pinned-original-source.tex, source lines 4457--4457. Cited by the retained statement or proof but not part of the collected spans; supplied only as a navigation pointer. The argument does not depend on its content.\par\medskip
\begin{proposition}
\label{proposition-stack-fpqc}
\begin{reference}
Proposition 3.3.6 of ``Intro to Algebraic Stacks'' by
Anatoly Preygel.
\end{reference}
Let $\mathcal{X}$ be an algebraic stack with quasi-affine\footnote{It suffices
to assume ind-quasi-affine.} diagonal. Then
$\mathcal{X}$ satisfies descent for fpqc coverings.
\end{proposition}

\begin{proof}
Our conventions are that $\mathcal{X}$ is a stack
in groupoids $p : \mathcal{X} \to (\Sch/S)_{fppf}$
over the category of schemes over a base scheme $S$ endowed
with the fppf topology. The statement means the following:
given an fpqc covering $\mathcal{U} = \{U_i \to U\}_{i \in I}$
of schemes over $S$ the functor
$$
\mathcal{X}_U \longrightarrow DD(\mathcal{U})
$$
is an equivalence. Here on the left we have the category
of objects of $\mathcal{X}$ over $U$ and on the right we have
the category of descent data in $\mathcal{X}$ relative to $\mathcal{U}$.
See discussion in Stacks, Section \ref{stacks-section-descent-data}.

\medskip\noindent
Fully faithfulness. Suppose we have two objects $x, y$ of $\mathcal{X}$
over $U$. Then $I = \mathit{Isom}(x, y)$ is an algebraic space over $U$.
Hence a collection of sections of $I$ over $U_i$ whose restrictions
to $U_i \times_U U_j$ agree, come from a unique section over $U$ by
the analogue of the proposition for algebraic spaces, see
Properties of Spaces, Proposition
\ref{spaces-properties-proposition-sheaf-fpqc}.
Thus our functor is fully faithful.

\medskip\noindent
Essential surjectivity. Here we are given objects
$x_i$ over $U_i$ and isomorphisms
$\varphi_{ij} : \text{pr}_0^*x_i \to \text{pr}_1^*x_j$
over $U_i \times_U U_j$ satisfying the cocyle condition
over $U_i \times_U U_j \times_U U_k$.

\medskip\noindent
Let $W$ be an affine scheme and let $W \to \mathcal{X}$
be a morphism. For each $i$ we can form
$$
W_i = U_i \times_{x_i, \mathcal{X}} W
$$
The projection $W_i \to U_i$ is quasi-affine as the
diagonal of $\mathcal{X}$ is quasi-affine. For each pair
$i, j \in I$ the isomorphism $\varphi_{ij}$ induces
an isomorphism
$$
W_i \times_U U_j =
(U_i \times_U U_j) \times_{x_i \circ \text{pr}_0, \mathcal{X}} W
\to
(U_i \times_U U_j) \times_{x_j \circ \text{pr}_1, \mathcal{X}} W =
U_i \times_U W_j
$$
Moreover, these isomorphisms satisfy the cocycle condition over
$U_i \times_U U_j \times_U U_k$. In other words, these isomorphisms
define a descent datum on the schemes $W_i/U_i$ relative to $\mathcal{U}$.
By Descent, Lemma \ref{descent-lemma-quasi-affine}
we see that this descent datum is effective\footnote{Or use
More on Groupoids, Lemma \ref{more-groupoids-lemma-ind-quasi-affine}
in the case of ind-quasi-affine diagonal.}.
We conclude that there exists a quasi-affine morphism
$W' \to U$ and a commutative diagram
$$
\xymatrix{
W' \ar[d] &
W_i \ar[l] \ar[d] \ar[r] &
W \ar[d] \\
U &
U_i \ar[l] \ar[r]^{x_i} &
\mathcal{X}
}
$$
whose squares are cartesian. Since $\{W_i \to W'\}_{i \in I}$
is the base change of $\mathcal{U}$ by $W' \to U$ we conclude
that it is an fpqc covering. Since $W$ satisfies the sheaf
condition for fpqc coverings, we obtain a unique morphism
$W' \to W$ such that $W_i \to W' \to W$ is the given morphism
$W_i \to W$. In other words, we have the commutative diagrams
$$
\xymatrix{
W_i \ar[d] \ar[r] &
W' \ar[d] \ar[r] &
W \ar[d] \\
U_i \ar[r] \ar@/_1pc/[rr]_{x_i} &
U &
\mathcal{X}
}
$$
compatible with the isomorphisms $\varphi_{ij}$ and whose
square and rectangle are cartesian.

\medskip\noindent
Choose a collection of affine schemes $W_\alpha$, $\alpha \in A$
and smooth morphisms $W_\alpha \to \mathcal{X}$ such that
$\coprod W_\alpha \to \mathcal{X}$ is surjective.
By the procedure of the preceding paragraph we produce
a diagram
$$
\xymatrix{
W_{\alpha, i} \ar[d] \ar[r] &
W_\alpha' \ar[d] \ar[r] &
W_\alpha \ar[d] \\
U_i \ar[r] \ar@/_1pc/[rr]_{x_i} &
U &
\mathcal{X}
}
$$
for each $\alpha$. Then the morphisms $W_\alpha' \to U$ are
smooth and jointly surjective.

\medskip\noindent
Denote $x_\alpha$ the object of $\mathcal{X}$ over $W_\alpha'$
corresponding to $W_\alpha' \to W_\alpha \to \mathcal{X}$.
Since $\mathcal{X}$ is an fppf stack and since
$\{W_\alpha' \to U\}$ is an fppf covering, it suffices
to show that there are isomorphisms
$\text{pr}_0^*x_\alpha \to \text{pr}_1^*x_\beta$
over $W_\alpha' \times_U W'_\beta$ satisfying the
cocycle condition. However, after pulling back to
$W_{\alpha, i}$ we do have such isomorphisms over
$W_{\alpha, i} \times_{U_i} W_{\beta, i} =
U_i \times_U (W_\alpha' \times_U W'_\beta)$
since the pullback of $x_\alpha$ to $W_{\alpha, i}$
is isomorphic to the pullback of $x_i$ to $W_{\alpha, i}$.
Since $\{U_i \times_U (W_\alpha' \times_U W'_\beta) \to
W_\alpha' \times_U W'_\beta\}_{i \in I}$
is an fpqc covering and by the aforementioned compatibility
of the diagrams above with $\varphi_{ij}$ these
isomorphisms descend to $W_\alpha' \times_U W'_\beta$
and the proof is complete.
\end{proof}

\end{document}
